\subsection{Activation Functions}\label{sec:activation-functions-in-cnns}

\subsubsection{Why CNNs Need Nonlinearities}

This section reconnects CNNs with something we already saw earlier for MLPs: \hl{stacking linear operations alone does not make the model nonlinear.}

\highspace
A convolutional layer is a linear/affine operation. If we flatten the input appropriately, it can be written in the familiar form:
\begin{equation*}
    \mathbf{a} = W\mathbf{x} + \mathbf{b}
\end{equation*}
The same is true for a dense layer (i.e., a fully connected layer). So if we \hl{stack only convolutional and dense layers}, without inserting nonlinear activations, the whole \hl{network can still collapse into one affine transformation.}

\highspace
For example, with two layers:
\begin{equation*}
    \mathbf{h} = W_1 \mathbf{x} + \mathbf{b}_1
\end{equation*}
And then the second layer:
\begin{equation*}
    \mathbf{y} = W_2 \mathbf{h} + \mathbf{b}_2
\end{equation*}
Substitute the first into the second:
\begin{equation*}
    \mathbf{y} = W_2 \left( W_1 \mathbf{x} + \mathbf{b}_1 \right) + \mathbf{b}_2 = W_2 W_1 \mathbf{x} + W_2 \mathbf{b}_1 + \mathbf{b}_2
\end{equation*}
Define:
\begin{equation*}
    W' = W_2 W_1, \quad \mathbf{b}' = W_2 \mathbf{b}_1 + \mathbf{b}_2
\end{equation*}
Then we have:
\begin{equation}
    \mathbf{y} = W' \mathbf{x} + \mathbf{b}'
\end{equation}
Which is still just one affine transformation. This is also the same argument used earlier in the image-classification section for MLPs (Multi-Layer Perceptron): \hl{multiple linear layers without nonlinearities are equivalent to a single linear classifier.}

\highspace
The same reasoning applies to CNNs because convolution itself is linear with respect to the input. So a stack such as:
\begin{equation*}
    \text{Conv}_{1} \rightarrow \text{Conv}_2 \rightarrow \text{Conv}_3 \rightarrow \dots
\end{equation*}
With no nonlinearities between them is \textbf{still equivalent}, in expressive power, to one larger linear transformation. That would defeat one of the main purposes of using a deep network: depth would not allow the model to represent genuinely nonlinear relationships.

\highspace
\textcolor{Green3}{\faIcon{check-circle}} \textbf{Activation functions solve this problem} by inserting a nonlinear transformation between linear layers:
\begin{equation*}
    \text{Conv}_{1} \rightarrow g(\cdot) \rightarrow \text{Conv}_2 \rightarrow g(\cdot) \rightarrow \text{Conv}_3 \rightarrow g(\cdot) \rightarrow \dots
\end{equation*}
For example:
\begin{equation*}
    \mathbf{h} = g\left(W_1 \mathbf{x} + \mathbf{b}_1\right)
\end{equation*}
Now, because $g(\cdot)$ is nonlinear, we cannot generally combine the two layers into a single matrix multiplication:
\begin{equation*}
    \mathbf{y} = W_2 g\left(W_1 \mathbf{x} + \mathbf{b}_1\right) + \mathbf{b}_2
\end{equation*}
Cannot be simplified to a single linear transformation. That is the fundamental reason \hl{activation functions increase the expressive power of the CNN.}

\highspace
However, it is important to note a fundamental detail: \hl{activation functions operate \textbf{element by element} on the activation volume.} Suppose a convolution produces:
\begin{equation*}
    A \in \mathbb{R}^{H \times W \times C}
\end{equation*}
An activation function $g(\cdot)$ produces:
\begin{equation*}
    A'(r, c, k) = g\left(A(r, c, k)\right)
\end{equation*}
So if:
\begin{equation*}
    A = \begin{bmatrix}
        -2 & 3 \\
        1 & -4
    \end{bmatrix}
\end{equation*}
And we layer use an activation function such as ReLU:
\begin{equation*}
    g(x) = \max(0, x) \quad \Rightarrow \quad A' = \begin{bmatrix}
        0 & 3 \\
        1 & 0
    \end{bmatrix}
\end{equation*}
\hl{Each number is transformed independently.} This also explains why activation layers normally \textbf{do not change tensor dimensions}. They change the \textbf{values}, not the shape.

\highspace
\begin{takeawaysbox}
    Convolutional layers are linear operations. Therefore, stacking several convolutional layers without activation functions would still yield an overall linear transformation and would not increase the expressive power of the network. Nonlinear activation functions are inserted between convolutional layers to allow the CNN to model nonlinear relationships. They operate independently on each activation value and therefore change the values of a feature volume without changing its shape.
\end{takeawaysbox}